\section{Kesimpulan}
Omni-RAG menyediakan konstruksi basis pengetahuan melalui \tech{chunker}, \tech{indexer}, dan pembangun \tech{knowledge graph} berbasis \tech{plugin}. Kontrak \tech{canonical artifact} mempertahankan identitas serta hubungan ke sumber sehingga strategi dapat diganti dengan batas integrasi yang sama.

Pada $K=25$, \tech{structure-aware chunker} mencapai \tech{recall} 0,740 pada \tech{User Manual} dan 0,64 pada peraturan. Untuk perbandingan keseluruhan dalam setiap korpus, \tech{dense} mencapai 0,835 pada $K=45$ untuk \tech{User Manual}, sedangkan \tech{hybrid} mencapai 0,67 pada $K=25$ untuk peraturan. Pada pertanyaan relasi hukum, \tech{sparse} mencapai \tech{KG recall} tertinggi sebesar 0,517 pada $K=25$. Hasil ini menunjukkan perlunya menyesuaikan strategi dengan struktur dokumen dan karakteristik pertanyaan.

Pengembangan berikutnya mencakup pengujian kombinasi komponen secara \tech{full factorial}, perluasan korpus dan pertanyaan pengguna, serta evaluasi jawaban akhir. Pengemasan \tech{plugin} juga perlu dilengkapi validasi versi kontrak, dependensi, dan pengujian kompatibilitas untuk mendukung penggunaan pada sumber baru.
